%% LyX 2.1.4 created this file.  For more info, see http://www.lyx.org/.
%% Do not edit unless you really know what you are doing.
\documentclass[a4paper,ngerman,naustrian,DIV=12,BCOR=1cm]{scrbook}
\usepackage[T1]{fontenc}
\usepackage[utf8]{inputenc}
\usepackage{fancyhdr}
\pagestyle{fancy}
\setcounter{secnumdepth}{3}
\usepackage{babel}
\usepackage{textcomp}
\usepackage{url}
\usepackage{makeidx}
\makeindex
\usepackage{graphicx}
\PassOptionsToPackage{normalem}{ulem}
\usepackage{ulem}
\usepackage[unicode=true,
 bookmarks=true,bookmarksnumbered=false,bookmarksopen=false,
 breaklinks=true,pdfborder={0 0 0},backref=false,colorlinks=false]
 {hyperref}
\hypersetup{pdftitle={Diplomarbeit Titel},
 pdfauthor={Wer auch immer},
 pdfsubject={Diplomarbeit},
 pdfkeywords={dies, das}}

\makeatletter

%%%%%%%%%%%%%%%%%%%%%%%%%%%%%% LyX specific LaTeX commands.
\pdfpageheight\paperheight
\pdfpagewidth\paperwidth

%% Because html converters don't know tabularnewline
\providecommand{\tabularnewline}{\\}

%%%%%%%%%%%%%%%%%%%%%%%%%%%%%% Textclass specific LaTeX commands.
\newcommand{\strong}[1]{\textbf{#1}}
\newcommand{\code}[1]{\texttt{#1}}

%%%%%%%%%%%%%%%%%%%%%%%%%%%%%% User specified LaTeX commands.
%%%%%%%%%%%%
% Latex-Vorspann
\usepackage{lastpage}
\usepackage{listings}
\usepackage{blindtext}

%% geht nicht mit jeder Latex Variante, gibt aber ein schöneres Layout
\usepackage{microtype}

%% Aufzählungen nicht so weit einrücken
\usepackage{enumitem}
%\setitemize{leftmargin=*}

%\usepackage{caladea}
%\usepackage[T1]{fontenc}
\usepackage{lmodern}

\usepackage{xspace}

%% für pandoc
%% maximale Breite der Bilder
\usepackage{graphicx}
\setkeys{Gin}{width=0.90\linewidth,keepaspectratio}
\providecommand{\tightlist}{%
\setlength{\itemsep}{0pt}\setlength{\parskip}{0pt}}

%%für Bilder nebeneinander
\usepackage{subfigure}

%um Schusterjungen und Hurenkinder vermeiden zu können
\usepackage{needspace}

%Farbpaket
\usepackage{xcolor}

%Für Bilddarstellung
\usepackage{float}

\makeatother

\usepackage{listings}
\addto\captionsnaustrian{\renewcommand{\lstlistingname}{\inputencoding{latin9}Listing}}
\addto\captionsngerman{\renewcommand{\lstlistingname}{\inputencoding{latin9}Listing}}
\renewcommand{\lstlistingname}{\inputencoding{latin9}Listing}

\begin{document}
%%%%%%
% Weitere Einstellungen siehe Latex-Vorspann

\sloppy % weniger Meldungen

\voffset5mm % etwas nach unten%%%%%%%%%%%%%%%%%%%%%%%%%%%%%%%%%%%%%%%%%%%%%%%%%%%%%%%%%%%%%%%%%%%%%%%%%%%%%%%%%%
% falls man die erste Zeile der Absätze nicht einrücken will
% dann sollte man aber etwas mehr Abstand zwischen den Absätzen erlauben
% Alternative \usepackage{parskip}
%%\setlength{\parindent}{0pt}
%%\setlength{\parskip}{1.5ex plus0.5ex minus0.5ex}
% Auch Fußnoten bündig ausrichten
\deffootnote[]{1em}{1em}{\textsuperscript{\thefootnotemark\ }}
% Listen etwas wenige einrücken, erfordert enumitem
\setitemize{leftmargin=*}

%%%%%%%%%%%%%%%%%%%%%%%%%%%%%%%%%%%%%%%%%%%%%%%%%%%%%%%%%%%%%%%%%%%%%%%%%%%%%%%%%%
%  Kopf und Fußzeilen -- links und rechts verschieden
\newcommand{\kopfseitenummer}{{\bfseries \thepage}}
\newcommand{\kopfkapl}{{\bfseries\leftmark}}
\newcommand{\kopfkapr}{{\bfseries\rightmark}}
\newcommand{\kopfbild}{\includegraphics[width=25mm]{Bilder/HTL3RLogoRGB}}
\newcommand{\kopfHTL}{Höhere Technische Bundeslehranstalt Wien 3, \\Rennweg 	Abteilung für Informationstechnologie}
\renewcommand{\chaptermark}[1]%
  {\thispagestyle{fancy}\markboth{\thechapter.\ #1}{}}%\thispagestyle{fancy}

%\lhead[\fancyplain{\kopfbild}{\kopfbild}]% li aussen
%      {\fancyplain{\kopfHTL}{\kopfHTL}}% re innen
%\rhead[\kopfHTL]% li innen
%      {\kopfbild}% re aussen

%% mit kapitelautor kann man den Autor festlegen oder auf leer setzen - steht dann in der Fußzeile.
\newcommand{\kapitelautor}{}

%%%
% Alternative: am Rand (Marginale)
%\setlength{\marginparsep}{-5mm}
%\mbox{}\marginpar{\raggedleft\hspace{0pt}Autor: Hans Huber}

%% kopf links: [linke] und {rechte} Seite
\lhead[\kopfbild]{\kopfkapl}
\rhead[\kopfkapr]{\kopfbild}
\chead{}

\lfoot[\kopfseitenummer]{\kapitelautor}
\cfoot[]{}
\rfoot[\kapitelautor]{\kopfseitenummer}
\renewcommand{\footrulewidth}{0.2pt}
\renewcommand{\headrulewidth}{0.2pt}

%%
% einfaches "siehe ..." - das Ziel muss man markieren
\newcommand{\kap}[1]{Kapitel~\ref{#1}, Seite~\pageref{#1}}
\newcommand{\siehe}[1]{siehe \kap{#1}}

%% http://ieg.ifs.tuwien.ac.at/~aigner/download/tuwien.sty
%Div. Abkürzungen (in Anlehnung an Jochen Köpper, jkthesis):
%\RequirePackage{xspace}
\newcommand{\bzw}{bzw.\@\xspace}
\newcommand{\bzgl}{bzgl.\@\xspace}
\newcommand{\ca}{ca.\@\xspace}
\newcommand{\dah}{d.\thinspace{}h.\@\xspace}
\newcommand{\Dah}{D.\thinspace{}h.\@\xspace}
\newcommand{\ds}{d.\thinspace{}s.\@\xspace}
\newcommand{\evtl}{evtl.\@\xspace}
\newcommand{\ua}{u.\thinspace{}a.\@\xspace}
\newcommand{\Ua}{U.\thinspace{}a.\@\xspace}
\newcommand{\usw}{usw.\@\xspace}
\newcommand{\va}{v.\thinspace{}a.\@\xspace}
\newcommand{\vgl}{vgl.\@\xspace}
\newcommand{\zB}{z.\thinspace{}B.\@\xspace}
\newcommand{\ZB}{Zum Beispiel\xspace}

%%%%%Programmlisting%%%%%%%%
% das braucht man nur einmal
% das kann auch ganz oben stehen
\lstset{basicstyle = \ttfamily\footnotesize,
	keywordstyle = \color{blue},
	stringstyle = \color{cyan},
	commentstyle = \color{lightgray},
	numbers = left,
	numberstyle = \tiny,
	stepnumber = 2,
	numbersep = 5pt,
	showspaces = false,
	frame = single,
	escapechar = $
}
%lstset{language=...} vor jedem Listing


% einmal oder immer was anderes
%vor jedem Listing: ********	\lstset{language=C}


%%%%%Anfang Titelseite
\pagenumbering{roman}
\title{Diplomarbeit}
\begin{titlepage}
\begin{minipage}[b]{1\columnwidth}
\parbox[b]{50mm}{\includegraphics[width=45mm]{Bilder/HTL3RLogoRGB}}
\hfill
\parbox[b]{130mm}{\footnotesize \textsc{Höhere Technische Bundeslehranstalt} Wien 3, Rennweg\\
IT \& Mechatronik\\
\\
HTL Rennweg :: Rennweg 89b\\
A-1030 Wien :: Tel +43 1 24215-10 :: Fax DW 18
}\\
\mbox{}
\end{minipage}

\vspace{1cm}


\begin{center}
\textbf{\LARGE{}Diplomarbeit}{\large{}}\\
{\large{}\vspace{15mm}
 }\textbf{\large{}}\\
\textbf{\large{}Hovering Steward}\\
 \vspace{15mm}
 ausgeführt an der\\
 Höheren Abteilung für Informationstechnologie/Ausbildungsschwerpunkt\\
 der Höheren Technischen Lehranstalt Wien 3 Rennweg\\
 \vspace{1cm}
 im Schuljahr 2015/2016\\
 \vspace{1cm}
 durch\\
 \vspace{0.5cm}
\textbf{\large{}Christina Bornberg}\\
\textbf{\large{}Katharina Joksch}\\
\textbf{\large{}Markus Kaiser}\\
\textbf{\large{}Alexander Punz}\\
\textbf{\large{}Lucas Ullrich}\\

\par\end{center}{\large \par}

\begin{center}
\vspace{20mm}
\normalsize unter der Anleitung von\\
\vspace{0.5cm}
Mag. Andreas Fink\\
DI Herbert Fleck
\par\end{center}

\begin{center}
\vspace{5mm}
Wien, \today
\par\end{center}

\end{titlepage}%%%%%%%%%%%%%%%%%%%%% Ende Titelseite %%%%%%%%%%%%%%%%%%%%%%


\chapter*{Kurzfassung}

% Auf Seiten mit einem neuen Kapitel ist keine Kopfzeile -- kann man sich aber wünschen
\thispagestyle{fancy}

Darum geht es.

\blindtext[1]


\chapter*{Abstract}

% mit Kopfzeile
\thispagestyle{fancy}

Thats why -- the translated text ,,Kurzfassung`` (this should be
an exact translation).

\blindtext[1]


\chapter*{Ehrenwörtliche Erklärung}

% mit Kopfzeile
\thispagestyle{fancy}

Ich versichere,
\begin{itemize}
\item dass ich meinen Anteil an dieser Diplomarbeit selbstständig verfasst
habe,
\item dass ich keine anderen als die angegebenen Quellen und Hilfsmittel
benutzt habe
\item und mich auch sonst keiner unerlaubten Hilfe bzw. Hilfsmittel bedient
habe.
\end{itemize}
\bigskip{}
Wien, am \today

<eigenhändige Unterschriften aller Teammitglieder>


\chapter*{Präambel}

\thispagestyle{fancy}

Die Inhalte dieser Diplomarbeit entsprechen den Qualitätsnormen für
,,Ingenieurprojekte`` gemäß §\,29 der Verordnung des Bundesministers
für Unterricht und kulturelle Angelegenheiten über die Reife- und
Diplomprüfung in den berufsbildenden höheren Schulen, BGBl. Nr. 847/1992,
in der Fassung der Verordnungen BGBl. Nr. 269/1993, Nr. 467/1996 und
BGBl. II Nr. 123/97.

\vspace{10mm}


\noindent Liste der betreuenden Lehrer:

<{[}Dir|AV|Prof{]}, akad. Grad, Vorname Name Hauptbetreuer>

<{[}Dir|AV|Prof{]}, akad. Grad, Vorname Name Hauptbetreuer Stellvertreter>

<{[}Dir|AV|Prof{]}, akad. Grad, Vorname Name Betreuer> ... (in alphabetischer
Reihenfolge des Nachnamens)

<{[}Dir|AV|Prof{]}, akad. Grad, Vorname Name Betreuer>

\vspace{10mm}


\noindent Liste der Kooperationspartner:%falls vorhanden

%%%%%%%%%%%%%%%%%%%%%%%%%%%%%%%%%%%%%%%%%%%%%%%%%%%%%%%%%%%%%%%%%%%%%%%%%%%%%%%%%%%%%%%%
%Verzeichnisse -- machen wir mit fancy headern
\renewcommand*{\chapterpagestyle}{fancy}
\cleardoublepage{}
\tableofcontents{}
\cleardoublepage{}
\listoftables
\cleardoublepage{}
\listoffigures

\cleardoublepage{}

%hier geht es los mit dem Text - auf einer rechten Seite
\pagenumbering{arabic}
\pagestyle{fancy}
\thispagestyle{fancy}


% wer hat diese Kapitel geschrieben oder leer
%\renewcommand{\kapitelautor}{Autor: Hans Huber}

\thispagestyle{fancy}

%%%%%%%%%%%%%%%%Inhalt der Diplomarbeit ab hier mit \include{}%%%%%%%%%%%%%%%%%%%%%%%%%%
%%Bilder werden meißt am besten mit [H] anstelle von [tbh] angezeigt


%%%%%%%%%%%%%%%%%%%%%%%%%%%%%%%%%%%%%%%%%%%%%%%%%%%%%%%%%%%%%%%%%%%%%%%%%%%%%%%%%%%%%%%%%%
% wer hat diese Kapitel geschrieben oder leer
\renewcommand{\kapitelautor}{}

\appendix

\chapter{Anhang 1\label{chap:Anhang-1}}

\printindex{}

\bibliographystyle{plaindin}
\bibliography{diplom}

\end{document}
